%% kapa: exam=2 round=37 year=2026 period=1 subject=appraisal_practice subject_ko=감정평가실무 seq=3 number=3 points=20
%% source: Q-Net 감정평가사 기출문제 (cst00309) artlSeq=5268128 file=raw/1008/5268128/26년 제37회 감정평가사 2차 (1교시_감정평가실무).pdf
%% conversion: kordoc 4.24.51 -> tools/gen2.py (auto); review=proofread-sonnet; flags=none
\begin{longquestion}{【문제 3】}
감정평가사 甲은 과수원(배나무 식재)의 매각 예정금액을 결정하기 위한 일반거래 목적의 감정평가를 의뢰받았다. 아래의 수집된 자료를 바탕으로 과수원의 감정평가액을 다음 물음에 따라 산정하시오. (20점) \par
\medskip\noindent 물음 1) 과수원 부지와 과수목이 일체로 거래된 사례가 없는 경우, 과수원 부지 및 과수목의 감정평가 방침을 기술하고, 과수원 부지(토지)만의 감정평가액을 산정하시오. (5점) \par
\medskip\noindent 물음 2) 감정평가 방침에 따른 과수목의 감정평가액을 산정하되, 수익수 및 유령수로 구분하여 산정한 후, 합산하시오. (15점) \par
< 자료 1 > 평가대상 자료 등 \par
1. 소재지: ○○시 ○○동 223번지 \par
2. 지목 등: ‘과수원’이고, 8m 포장도로에 접한 사다리형의 평지 \par
3. 면적: 5,000m\textsuperscript{2} \par
4. 현황: 과수원이며, 15년생 배나무(수익수) 300주, 3년생 배나무(유령수) 100주, 정상식재 \par
5. 용도지역: 계획관리지역 \par
6. 기준시점: 2026. 1. 1. \par
< 자료 2 > 토지관련 가격자료 \par
1. 인근지역의 표준지 공시지가: 2026. 1. 1. 공시된 인근지역의 표준지는 본건과 제반 가치형성요인이 유사(지목, 현황, 용도지역 등)하며, 공시지가는 40,000원/m\textsuperscript{2}임 \par
2. 지가변동률: 중소도시 외곽 순수 농경지대로 변동사항 없음 (0\%) \par
3. 거래사례: 최근 3년 동안 인근지역과 동일수급권 내의 유사지역을 포함하여 과수원과 과수목을 일체로 한 거래사례는 포착할 수 없었고, 최근 배의 수출 증가로 농경지를 과수원으로 전환하려는 수요로 인한 거래가 포착되며, ‘전’과 ‘답’, ‘과수원’ 등의 구분없이 공시지가 수준에서 거래가격이 형성되는 것으로 조사됨 \par
< 자료 3 > 수익자료 \par
1. 15년생 배나무는 현재 연간 주당 60,000원의 순수익을 창출하며, 이러한 수익은 향후 5년간 지속된 후, 모두 수종 갱신될 예정이고, 매년 순수익 증가율은 2\%로 예상됨 \par
2. 상기 순수익은 통상의 관리비용이 공제된 후의 금액이고, 할인율은 농업부문의 위험률 특성을 반영한 연간 10\%임 \par
< 자료 4 > 원가자료 \par
1. 정상수확기에 미달한 3년생 배나무는 묘목상태로 3년전 주당 10,000원에 식재되었고, 현재의 묘목포함된 식재비용도 주당 10,000원 내외로 조사되어 취득원가에 대한 이자는 고려하지 아니함 \par
2. 유령수만의 관리비용은 연간 평균기준으로 주당 5,000원 수준으로 파악되며, 금액적 중요성이 크지 아니하다고 판단되어 이의 투입시기에 따른 이자는 고려하지 아니함 \par
3. 유령수인 3년생 배나무의 현재 상태로의 가격은 성장에 따른 이식에서의 인건비 과다로 경제성이 없어 별도의 판매는 어려운 상황임 \par
\end{longquestion}
